\section{Expanded reasoning: page-to-line segmentation}
\label{sec:expand-segment}
\subsection{The measurement model and why normalise lighting}
The input is a grey image $G(y,x)$, with row $y$ increasing downwards and
column $x$ increasing rightwards. A pixel is an area sample, not a physical
point. A useful approximate model is $G=L\,R_s+\eta$: $L$ is slowly varying
illumination, $R_s$ surface reflectance and $\eta$ sensor/readout noise.
Paper and ink have different reflectance. The method estimates the local
paper level $B_r$ and divides by it. If $B_r\simeq L R_{\rm paper}$,
then $G/B_r\simeq R_s/R_{\rm paper}$ and the multiplicative illumination
approximately cancels. This algebra is the reason for using a ratio rather
than a fixed grey threshold; it is not proof that strong shadows or glare
are removed. Additive glare, saturation, texture and a window mostly filled
with ink violate the approximation. Sauvola's local mean/variance threshold
is a published alternative~\cite{sauvola2000}; choosing it would change the
measurement model and require independent comparison.

\paragraph{Window size and edges.}
A window must be wider than an individual stroke so that it estimates paper
rather than that stroke, yet narrower than the illumination variation it
is supposed to track. Too small a window follows and cancels ink; too large
a window misses local shadows. Edge replication extends the nearest edge
pixel outside the image. Zero padding would invent black borders. Reflection
and reduced edge windows are alternatives. The selected radius is an
implementation setting, not a universal imaging constant.

\subsection{Integral image: four lookups, with an arithmetic example}
For a zero-padded prefix-sum array define
\[
S(y,x)=\sum_{0\leq i<y}\sum_{0\leq j<x}G(i,j),\qquad S(0,x)=S(y,0)=0.
\]
The half-open rectangle $[y_0,y_1)\times[x_0,x_1)$ has sum
\[
Q=S(y_1,x_1)-S(y_0,x_1)-S(y_1,x_0)+S(y_0,x_0).
\]
The two strips are subtracted, and the doubly subtracted intersection is
added back. This is inclusion--exclusion. Crow introduced summed-area tables
for efficient area aggregation~\cite{crow1984}. Constructing $S$ takes
$O(HW)$ additions and memory; each subsequent rectangle sum takes constant
work regardless of radius. The complete image filter still takes $O(HW)$.
It does not mean the whole operation takes constant time.

For $G=\left[\begin{smallmatrix}1&2&3\\4&5&6\\7&8&9\end{smallmatrix}\right]$,
the padded $S$ is
$\left[\begin{smallmatrix}0&0&0&0\\0&1&3&6\\0&5&12&21\\0&12&27&45\end{smallmatrix}\right]$.
The lower-right $2\times2$ sum is $45-6-12+1=28$, equal to $5+6+8+9$.
Its mean is $28/4=7$. Rectangle boundaries and padding must use the same
convention: an off-by-one error changes the area and its denominator.
Prefix sums should accumulate in a type large enough to avoid overflow.

\subsection{Ratio, ink/paper distinction and Otsu's criterion}
The report's ratio is $R=G/\max(B_r,1)$, with grey levels on the 0--255
scale. The floor of 1 is a numerical guard measured in grey-level units;
it is not a claim that a near-black local background is reliable. For the
reported example $G=100$, $B_r=161.43$, $R\simeq0.6195<0.80$, so the pixel
is ink~\cite{finalreport2026}. The brightness ratio is different from an
\emph{ink-area fraction} $\sum M/(HW)$, which counts classified pixels.

For histogram probabilities $p_i=n_i/N$, split at bin $t$. Set
\[
\omega_0=\sum_{i\le t}p_i,\quad \omega_1=1-\omega_0,\quad
\mu_0={\sum_{i\le t}i p_i\over\omega_0},\quad
\mu_1={\sum_{i>t}i p_i\over\omega_1}.
\]
Otsu selects $t$ maximising $\omega_0\omega_1(\mu_0-\mu_1)^2$.
This is between-class variance; total variance equals within-class plus
between-class variance, so maximising one minimises the other for a fixed
histogram~\cite{otsu1979}. Empty classes are excluded. ``Optimal'' here
means optimal for that histogram objective, not for true character recovery.
The project bins ratios into 256 bins over 0--1.2, giving bin width
$1.2/256=0.0046875$. Bin count and range are numerical resolution choices.

Constructed example: four samples at $0.4$ and six at $1.0$ have
$\omega_0=0.4$, $\omega_1=0.6$, and between-class variance
$0.4(0.6)(0.4-1)^2=0.0864$. Any threshold strictly between the two levels
gives the same partition. Thus an Otsu threshold need not be unique.
On a unimodal almost-blank image it can separate paper noise rather than ink.
Inspecting the resulting mask remains necessary.

\subsection{Orientation, deskew, run limits and the line-pitch estimate}
A quarter-turn and a small-angle deskew solve different problems. A 90-degree
turn changes which image axis follows the lines. A fine rotation corrects a
small slope. The report requires the nearly-empty-row score to improve by
a factor 1.35 before accepting the quarter-turn. If scores are $0.20$ and
$0.28$, the ratio is 1.4 and passes; $0.20$ and $0.25$ do not. The margin
expresses conservatism when evidence is weak. Neither 1.35 nor this score
distinguishes upright text from upside-down text in general.

For each trial angle $\theta$, the row projection is
$p_\theta(y)=\sum_x M_\theta(y,x)$ and the objective is
$E(\theta)=\sum_y p_\theta(y)^2$~\cite{postl1986}.
For fixed total ink, concentrating it gives a larger squared sum:
$(8,0)$ scores 64 whereas $(4,4)$ scores 32. Horizontal rows therefore tend
to win. The 17-point grid from $-4$ to $4$ degrees in half-degree steps
is a search range and resolution choice. If the true optimum lies inside
the grid and the objective is well behaved, nearest-grid angular error is
at most 0.25 degrees. That does not bound estimation error on curved text.
Rotating sparse coordinates during the search saves resampling; the chosen
full grey image must then be rotated and binarised consistently.

The line pitch $P$ is a typical centre-to-centre distance in working-image
pixels. It is not x-height. Pitch-normalised guards adapt to image scale.
The report's horizontal and vertical run limits are $1.25P$ and $1.4P$;
at $P=80$ pixels these are 100 and 112 pixels. Long straight ruling is
more likely to exceed them than an ordinary stroke. Long underlines and
vertical letter strokes can nevertheless be removed. These are heuristics,
and estimating $P$ again after mark removal reduces circular dependence.

For a smoothed row profile, the 6\%-of-maximum candidate floor rejects weak
peaks. A maximum of 200 implies a candidate floor of 12, in the same units
as the profile. If one dark line dominates, faint genuine lines may fall
below the floor. Sorting peaks by strength and suppressing others within
$0.6P$ avoids treating one tall text body as multiple lines. At $P=80$,
the separation is 48 pixels. Increasing it merges tightly spaced real lines;
decreasing it admits duplicate centres. Median peak spacing resists isolated
bad gaps. Autocorrelation $A(k)=\sum_y p(y)p(y+k)$ instead seeks a repeated
spacing; strong ruling or harmonics can select the wrong multiple. Component
heights supply a fallback scale when peak evidence is sparse. None of these
fallbacks independently verifies the transcript's line count.

\subsection{Seam cost, closing and dynamic programming, fully unpacked}
The seam is one row $s_x$ for every column $x$, constrained to the corridor
between neighbouring line centres. For the report's cost
\[
C(y,x)=4M_c(y,x)+0.012\left({y-y_m\over h_g}\right)^2,
\]
$M_c$ is the closed binary mask; $y_m$ is the corridor midpoint and $h_g$
its half-height. Both terms are dimensionless. The quadratic prior mildly
favours the gap's middle; the ink term strongly discourages stroke crossings.
Within $|y-y_m|\le h_g$, the prior is at most 0.012 while one ink pixel
costs 4. This 333-to-1 local ratio explains priority, not an experimentally
proved optimum. Total route cost depends on every column, so a long detour
can outweigh one local penalty.

Closing is dilation followed by erosion with a chosen structuring element.
A square radius $r$ has side $2r+1$. The report uses
$r=\max(2,\lfloor0.07P\rfloor)$; $P=80$ gives $r=5$ and side 11.
It discourages a seam from slipping through tiny holes in a stroke. Too
large a radius joins neighbouring lines and destroys real passages. The
maximum here is a \emph{lower bound} of 2 on radius, not a maximum permitted
radius. The phrase ``max closing radius'' in the supplied list should not be
read as an upper cap unless a particular implementation adds one.

For allowed vertical changes $d\in\mathcal D$, accumulated cost is
\[
D(y,x)=C(y,x)+\min_{d\in\mathcal D}
       [D(y+d,x-1)+0.02|d|].
\]
The 0.02 term prices vertical movement and reduces jagged seams. Displacement
is measured in working-image rows. Restricting $\mathcal D$ constrains slope;
its exact contents belong to the implementation. Store each minimising
predecessor, choose the lowest-cost final state, and backtrack. Without the
predecessors, the minimum cost alone is not a path. This is Bellman's optimal
substructure principle on an acyclic graph~\cite{bellman1954,seampaperlink}.

Constructed arithmetic, for $\mathcal D=\{-1,0,1\}$: if predecessor costs
at rows $y-1,y,y+1$ are $0.10,0.16,0.11$, the transition candidates are
$0.12,0.16,0.13$. A blank pixel with prior 0.003 then has $D=0.123$.
An ink pixel at that location has $D=4.123$. Greedily choosing the cheapest
pixel in each column can fail because it ignores a later blocked corridor.
Dynamic programming compares complete prefixes. Its optimum is global only
within this corridor, step set and cost function, not over all conceivable
line separations. With $W$ columns, $H_g$ corridor rows and $q$ transitions,
time is $O(WH_gq)$; backtracking storage is $O(WH_g)$.

\subsection{Component ownership and the 15\% substantial-band rule}
Seam bands assign space, not semantic identity. Let $a_{C,n}$ be the ink
pixels of connected component $C$ inside band $n$. A substantial band obeys
$a_{C,n}\ge0.15\sum_j a_{C,j}$. With overlaps $(80,15,5)$ and total 100,
the first two qualify; the third does not. The topmost substantial band is
a reference owner for eligible corrections, not an instruction to transfer
the whole component blindly. With threshold 0.15, at most six disjoint bands
can qualify because seven would require at least 105\% of the component.
This bound follows from the definition; it does not justify choosing 0.15.

Eight-connectivity includes diagonal neighbours. Run-based labelling stores
contiguous horizontal ink intervals and unions runs touching the row above,
including diagonal overlap~\cite{he2008,tarjan1975}. It avoids recursively
visiting every pixel in Python. Raw components are revisited because closing
can create artificial connections. Dot-to-stem, ring and residual-fragment
rules use size/proximity guards to avoid moving a genuine detached feature
to another line. A dot's meaning cannot be recovered from proximity alone.
The report's 462 reassigned pixels describe one recorded example; they are
not 462 verified corrections and not an accuracy percentage.

\subsection{Contrast, darkness exponent and solid-crop rejection}
Contrast stretching maps chosen dark/light endpoints to a wider output range.
For darkness $D=1-G/255$ in $[0,1]$, the transformation
$G'=255(1-D^{0.55})$ strengthens intermediate darkness: $D=0.25$ becomes
approximately 0.467. Applying the same exponent directly to brightness
would have a different effect. Always identify which variable is exponentiated.
The coefficient 0.55 is a project contrast heuristic, not a gamma correction
derived from a known camera response.

The supplied list also asks about a dark-pixel fraction exceeding 15\% and
three times the page median. In the inspected \texttt{DropSolidCrops}, this
applies to whole candidate line crops, not individual pixel rows: darkness
is grey level below 128, and the median is taken across candidate crops.
For median
0.04, the operative threshold is $\max(0.15,0.12)=0.15$; for median 0.08
it is 0.24. Source comments report 668 crops with densest genuine crop 12.37\%,
largest genuine-to-page-median ratio 2.18, and one artefact at 37.5\% and
7.2 times its median. That supplies documented local intent, not independently
rerun validation or a literature optimum. If every crop would be rejected,
the inspected function returns the original list. These tests target unusually
dense artefacts and can also reject heavy handwriting. A factor of three is a relative guard, not a
three-standard-deviation statistical test. The 15\% in this test has a
different denominator and purpose from the 15\% ownership test.

\subsection{Complete algorithm, alternatives and verification meaning}
The conceptual order is: decode greyscale; reduce the working resolution;
estimate background and ink; choose quarter-turn and fine skew; estimate pitch;
remove eligible long marks; re-estimate pitch and centres; close a mask for
seam costs; compute and backtrack boundaries; assign initial bands; refine
ownership using raw and closed components; crop each line from the grey image;
exclude neighbouring ink; apply the chosen crop preparation; sort by page order.
The crop used for classification must preserve the mapping back to the original
pixels used for glyph extraction.

Straight projection cuts are simpler but cannot bend around ascenders;
connected-component clustering is natural for print but cursive joins and
detached dots complicate it; learned segmentation needs labelled boundaries
and another trained model. Seams offer explicit, inspectable geometry.
That is a defensible method trade-off, not proof of superiority on every page.

Verification asks whether sums, rotations, DP recurrence and ownership obey
their definitions. Validation asks whether the resulting lines contain the
intended handwriting under the photographed conditions. Equal crop counts
can conceal a missing line and a duplicated line. A useful seam plot shows
the raw mask, closed mask, cost scale, centre corridor, final path and ownership
colours. A path in a low-cost valley explains the optimisation; only comparison
with labelled ink establishes correct separation. No old-segmenter comparison,
new lighting test or sweep was performed for this edition.
